\makeabstract
    {
    Removing Barriers to Alcohol Use Disorder Treatment: Public Interest in Over-the-Counter Naltrexone
    }
    {
    Bardetti, L.\textsuperscript{1,3}, Rosenblum, N.\textsuperscript{1}, Engel, S.\textsuperscript{1}, Campbell, D.E.\textsuperscript{1}, Garcia, M.\textsuperscript{1}, Abbasi, S.\textsuperscript{1},  Whee, A.\textsuperscript{1}, Vatsyayan, A.\textsuperscript{1}, Suzuki, J.\textsuperscript{1,2}, Terechin, O.\textsuperscript{1,2}
    }
    {
    \textsuperscript{1} Department of Psychiatry, Mass General Brigham, Boston, MA\\
\textsuperscript{2} Harvard Medical School, Boston, MA\\
\textsuperscript{3} Amherst College, Amherst, MA
    }
    {
    This study aims to explore public interest in over-the-counter (OTC) naltrexone and assess whether this approach may serve as a feasible strategy to expand access to alcohol use disorder treatment. Individuals{\textquoteright} awareness, attitudes, and concerns toward OTC naltrexone have not yet been systematically studied, so this information is needed to guide future policy and education efforts to increase access to safe, low-barrier, and effective treatment for AUD.
    }
    {
    An over-the-phone survey assessed alcohol use history, awareness of naltrexone, willingness to use naltrexone if it was available OTC, perceived benefits and concerns related to OTC use, and preferred sources of guidance if such a medication were available without a prescription. Inclusion criteria were English-speaking, 18+ adults with self-reported alcohol use at any level.
    }
    {
    The results of this retrospective study indicate statistically significant positive effects of mammoth stuffed animals on student well-being. Analysis of the quantitative data from surveys revealed a significant reduction in students' reported stress levels (p \ensuremath{<} 0.001) and a notable increase in self-esteem (p \ensuremath{<} 0.05) when interacting with these plush companions. Moreover, qualitative insights obtained from interviews and observations highlighted the presence of mammoth stuffed animals as a significant factor contributing to a sense of security and emotional support among students. The findings emphasize the potential efficacy of unconventional interventions in promoting student welfare and underscore the importance of considering innovative approaches to enhance well-being in educational settings.
    }
    {
    Targeted Demand: Individuals diagnosed with AUD (84.3\% aware, 65.7\% likely/highly likely to use) and past naltrexone users (67.3\% likely/highly likely to use, 90.4\% likely/highly likely to recommend) showed the greatest interest in OTC access. Broad Acceptance: Across AUDIT-C groups and familiarity levels, over 80\% anticipate a positive impact from OTC naltrexone and would recommend to others. More education needed: Concerns highlight the need for clear instructions due to lack of medical supervision in an OTC setting. In conclusion, this retrospective study provides compelling evidence that mammoth stuffed animals have a positive impact on student well-being. The presence of these endearing companions is associated with reduced stress levels, increased self-esteem, and a sense of emotional security. These findings support the integration of unconventional interventions to enhance student welfare in educational settings.
    }

    \newpage
    